\scanpage{5}

Par suite, on a
\[
\bigwedge_{\underline O_X}^p J/J^2\simeq
f^*\Bigl(\bigotimes_i\bigwedge^{\alpha_i}M_{d_i}\Bigr)(-\sum\alpha_i d_i)
\]
où $M_{d_i}$ sont les composantes non nulles de $M$, supposées
de \unsure{degré} $\alpha_i$. D'où un isom
\[
\Omega^n_{X/S}\otimes f^*\Bigl(\bigotimes_i\bigwedge^{\alpha}M_{d_i}\Bigr)(-\sum\alpha_i d_i)
\xrightarrow{\sim}f^*\Bigl(\bigwedge_{\unsure{S|\underline O_S}}^{r+1}\underline E\Bigr)(-r-1)
\]
[Comprendre : un hom qui existe en tous cas dès
qu'on se donne des $M_{d_i}\to\operatorname{Symm}^{d_i}(\underline E)$, les $M_{d_i}$ loc.
libres de rang $\alpha_i$, et $X$ défini par
\inserted{l'idéal engendré par} les images
des $M_{d_i}$, posant $p=\sum\alpha_i$, $n=r-p$], i.e. on a
\[
\Omega^n_{X/S}\xrightarrow{\sim}
f^*\Bigl(\bigl(\bigwedge_S^{r+1}E\bigr)\otimes
\bigl(\bigotimes_i\bigwedge^{\alpha_i}\check M_{d_i}\bigr)\Bigr)
(\sum\alpha_i d_i-r-1)
\]

Dans le cas où
\[
\sum\alpha_i d_i=r+1
\]
on trouve donc
\[
\Omega^n_{X/S}\xrightarrow{\sim}f^*(-)
\]

Supposons p.ex. que $E$ soit le faisceau loc. libre de
rang $3$ associé à une courbe elliptique lisse
sur $S$, donc $r=2$, il n'y a qu'un seul $M_{d_i}$, $d_i=3$,
$\alpha_i=1$, on trouve
\[
\Omega^1_{X/S}\xrightarrow{\sim}f^*\bigl(\bigwedge^3E\otimes\check M\bigr)
\]
Or $E$ a une filtration de facteurs $\underline O_S$,
\deleted{$\unsure{g}$} $\mathcal L^{\otimes2}$, $\mathcal L^{\otimes3}$,
où $\mathcal L$ est inversible \unsure{tel que}
$f^*(\mathcal L)\simeq\unsure{\mathcal T}_{X/S}$, donc
$\bigwedge^3E\simeq\mathcal L^{\otimes5}$, et comme
$\Omega^1_{X/S}\simeq f^*(\check{\mathcal L})$, on aura
\[
\check{\mathcal L}\xrightarrow{\sim}\mathcal L^{\otimes5}\otimes\check M
\qquad\text{d'où}\qquad M\xrightarrow{\sim}\mathcal L^{\otimes6}.
\]

\mspage{4}
